\section{Research status, verification, and next questions}
\label{sec:audit}

\subsection{What has been established}

The proof-bearing deliverable is the article itself. The programs
check finite algebra and independent numerical instances; they are
not substituted for its convergence, continuation, contour, or
asymptotic arguments. The principal advances can be mapped precisely
to the source questions.

\begin{longtable}{@{}>{\raggedright\arraybackslash}p{0.24\textwidth}>{\raggedright\arraybackslash}p{0.45\textwidth}>{\raggedright\arraybackslash}p{0.25\textwidth}@{}}
\toprule
Source question & Result supplied here & Remaining scope\\
\midrule
\endfirsthead
\toprule
Source question & Result supplied here & Remaining scope\\
\midrule
\endhead
Individual ordered mixed labels &
Theorem~\ref{thm:ordered} gives the exact two-particle identity on an
open complex domain. Theorems~\ref{thm:poles} and
\ref{thm:digamma-master} give its continuation and all positive-order
values. &
Higher rank is handled by finite localization, not a universal
single-zeta assertion.\\[5pt]
Noninteger cutoff generators &
Theorem~\ref{thm:cutoff} and \eqref{eq:b-beta-lerch} give a finite
polylogarithm and beta--Lerch transform, entire in both orders on
the disk. &
The endpoint requires its own convergence domain and collective
cancellations.\\[5pt]
Other annihilating stencils &
Section~\ref{sec:localization} proves a finite block formula for
critical integer stencils and a stronger reduction under one further
annihilation condition. &
This is a sufficient family, not a classification of arbitrary
noninteger grids.\\[5pt]
All unequal even harmonic-tail prefixes &
Section~\ref{gb:section} proves eventual optimality on every fixed-gap
ray, with exact Gamma normalization and a nonzero signed asymptotic. &
No effective threshold is provided. The full all-pairs conjecture,
including complete gap-six coverage, remains open here.\\[5pt]
Stieltjes derivatives and primitives &
Theorems~\ref{thm:subtracted-jets} and \ref{thm:beta-stieltjes}
provide ordinary all-order identities. Equations
\eqref{eq:gamma-primitive} and \eqref{eq:arcsine-primitive}
give compatible primitives. &
Elementary symmetric arcsine cases are recurrence consequences;
they are not independent priority claims.\\[5pt]
Short Gaussian $S_6$ and revised $S_8$ &
Their present source status has been checked and preserved. &
Neither identity is proved or refuted by this article.\\
\bottomrule
\end{longtable}

The incoming bilinear and tail-kernel reports already supply complete
smaller-order-ten and smaller-order-twelve reconstruction
theorems \cite{repo:bilinear,repo:tail-kernels}. They should
retain their existing attribution and audit status. Reproving or
relabeling an already supplied finite case would not constitute a
new research advance; the fixed-gap theorem instead changes which
infinite direction is controlled.

\subsection{Audit findings and a proposed correction}

No newly false inherited mathematical theorem was identified in the
material used for this article. There is a concrete editorial defect:
the manuscript README repeats consecutively the complete paragraph
describing the 1,032-input first-letter-$X$ family, its rank 1,008,
and its kernel dimension 24 \cite{repo:manuscript}. The package includes
a narrowly scoped patch removing the second copy. This is a proposed
README correction, not a modification of a mathematical statement.

Several distinctions matter when the new results are integrated.
The source already records that the old $S_8$ candidate was rejected
and the revised candidate is still conjectural. The small residual
enclosures for $S_6$ and revised $S_8$ establish proximity, not zero.
Likewise, the canonical modular relation searches do not assert
rational nonmembership or a counterexample to a period equality
\cite{repo:ledger,repo:gaussian}. These are inherited status safeguards,
not newly discovered errors.

For the present contribution, the analogous precise statements are:
the zero-weight limits in Section~\ref{sec:continuation} concern a
meromorphic continuation; the constant in \eqref{eq:hard-cutoff} is
not the fixed-coordinate Laurent constant; block formulas outside
their separated convergence domain are interpreted collectively;
and full rank of a formal moment matrix does not prove arithmetic
independence of the evaluated constants. The Gamma argument also
shows why positivity of the scalar recurrence does not imply
positivity of every minor: even modes lose their leading Laplace
phase, and the determinant has the parity-dependent sign derived
in Section~\ref{gb:section}.

\subsection{Reproducibility and the independent checks}

The package contains modular TeX, a self-contained assembled TeX
file, the compiled PDF, a build command, a claim-status file, exact
source locations, and portable verification programs. Source
repository material was inspected at the pinned revision and kept
unchanged. The generated correction patch is separate.

The ordered-boundary suite checks the finite summation-by-parts
identity on 240 arbitrary rational sequence/cutoff cases, the rational
partial-fraction identity on 36 order pairs, 570 parity coefficients,
420 coefficients of the cutoff generator, and all four entries of
the zero-weight directional table. Four explicit infinite-sum values
are also compared with directly accumulated original ordered sums,
using both positive tail envelopes and separately labelled accelerated
diagnostics.

The corollary suite checks 13 exact master-function derivatives and
special values and nine exact formal Stieltjes coefficients. It
compares the geometric master with the original resolvent-weighted
ordered sum, tests two subtracted harmonic sums, follows six
unaccelerated hard-cutoff limits, and evaluates nine arcsine integrals.
The beta--Stieltjes suite checks five Bell--Leibniz coefficients,
independent Stieltjes-value samples, and twelve numerical beta
integrals through jet order two, including independently varied
shifts. Its expansion and endpoint truncations are recorded
separately from floating-point roundoff.

The localization suite supplies 147 exact checks, including
finite-support transport, the full contact table, both rank-three
two-index sectors, singleton cancellation through rank five, and a
deliberately corrupted control. The Gamma suite compares directly
integrated basis coefficients with the rational recurrence for
$m=1,2,3$, checks the exact inhomogeneous moment, and records
finite-index ratios to the determinant asymptotic. A separate exact
suite checks the Bernoulli matrix against its complementary-minor
normalization and the harmonic correction. The independent proof
audit rederives the contour shifts, both residues, the initial
determinant, the parity-sensitive asymptotics, and the source bridge.
This is internal independent checking, not external peer review or
proof-assistant verification.

All numerical comparisons are diagnostics at their stated working
precisions. The proof of eventual nonvanishing is the nonzero
limiting Vandermonde determinant, not an extrapolation of the
displayed finite data. In particular, a finite gap-six sign table
does not provide the missing effective cutoff.

\subsection{Further research questions}

The following questions are arranged by how directly the present
proofs supply a route forward. They concern exact identities and
structural classification; estimates enter only where needed to
turn an eventual identity theorem into a finite decision.

\begin{enumerate}[label=\textbf{\arabic*.},leftmargin=*]
\item \textbf{Complete the fixed-gap prefix conjectures.}
For each fixed $m$, turn the nonzero asymptotic into an explicit
computable threshold $B(m)$ and combine it with exact rational
determinants below that threshold. The specific next case is gap
six. The exact values at $b=4,5,6$ are negative, and the checked
values $7\le b\le40$ are positive. The strengthened conjecture
\[
 \det F_{b,3}>0\qquad(b\ge7)
\]
is consistent with both the finite data and the proved eventual
sign, but is not proved here. A successful completion must justify
the transition from a finite table to all indices.

\item \textbf{Evaluate the Gamma normalization constants.}
The positive integral $K_m$ is exact, and $K_1=8\pi$.
Find closed Gamma, hypergeometric, or modular descriptions of
$K_2,K_3,\ldots$. The multiple-orthogonality interpretation suggests
studying the associated moment determinant rather than integrating
each entry separately. A recurrence in $m$, if it exists, would be
particularly useful.

\item \textbf{Connect the two eventual regimes.}
The new theorem fixes the gap while the earlier theorem fixes the
smaller order. Can one obtain a two-parameter determinant expansion
that covers the region between those rays? The limiting nodes
$(j\pi)^{-2}$ become less separated as $m$ grows, so a proof must
control the full determinant, not merely reuse fixed-$m$ error
constants. This is a possible route toward the all-pairs conjecture.

\item \textbf{Classify the finite sign changes.}
The exact Gamma representation is valid at every $b>m$, while its
sign is presently asymptotic. Determine the number and locations of
the finite sign transitions for each $m$. A variation-diminishing
statement adapted to the polynomials
$\operatorname{Im}(x+i/2)^{2n}$ could replace a false appeal to
unqualified total positivity.

\item \textbf{Deform the Gamma basis.}
Replace the quarter shift by a parameter and seek the corresponding
shift equation, residue forcing, and reconstruction matrix. Identify
which deformations preserve a finite set of sinh/tanh modes and
which create additional residue terms. Differentiating a valid
deformation would give new exact polygamma-weighted orthogonality
identities; the deformation itself must be proved before its
derivatives are used.

\item \textbf{Determine the actual higher-rank period coordinates.}
The block theorem gives at most
$\lfloor(r+1)/2\rfloor$ free ordered indices. At rank three the
two-index contact sector is nonzero, but this does not prove
arithmetical depth two. Determine which integer orders, shift
collisions, or label symmetries reduce these sectors to single
Hurwitz values, and which require genuine nested periods.
The explicit contact polynomial provides a finite starting point.

\item \textbf{Compute all-rank polar geometry from the blocks.}
The two-particle oblique residues are explicit polynomials.
At higher rank, the finite sector formula should organize polar
hyperplanes by free blocks and their mergers. Prove a complete
local Laurent description, including its behavior when the original
finite-difference weights vanish. This would extend the directional
table without confusing a pointwise specialization with a limit.

\item \textbf{Classify stencils by ordered boundary support.}
The annihilation conditions proved here are sufficient. Seek
necessary conditions for a prescribed maximum number of free
blocks, first for finite integer stencils. A discrete-residue
criterion might recognize additional cancellations invisible to
degree counting. Nonuniform or noninteger grids require a separate
analysis because there is no common integer transport lattice.

\item \textbf{Find new exact beta--Stieltjes evaluations.}
For rational $\sigma$, special shifts may combine the beta weight
with multiplication or reflection formulas. Determine when
\eqref{eq:beta-stieltjes-all} reduces to a finite expression in
zeta derivatives, log-Gamma values, and lower Stieltjes constants,
rather than the generally infinite but geometric expansion
\eqref{eq:beta-jet-zeta}. Confluent shift derivatives offer a
systematic way to introduce polygamma factors.

\item \textbf{Construct higher-rank cutoff formulas.}
Insert a common outer cutoff before transporting the stencils.
Track the factors of $1-z$ attached to free collision blocks and
determine whether the endpoint depth reduction can be read directly
from them. Section~\ref{sec:polylog} proves this mechanism at rank
two. A higher-rank formula should retain all finite endpoint terms
before any limiting operation.

\item \textbf{Return to Gaussian periods with a new functional route.}
The short $S_6$ and revised $S_8$ questions remain worthwhile, but
the inspected finite relation spaces have already been analyzed
carefully. New work should supply a functional identity, a genuinely
larger motivic or iterated-integral relation, or a rigorously changed
candidate. Further decimal agreement alone cannot settle either
equality. The ordered kernels here may suggest new generating
relations, but no such implication is presently proved.
\end{enumerate}

\subsection{Suggested integration}

The ordered identity and its cutoff generator naturally continue
the finite-difference report. The beta--Stieltjes identities belong
with the Hurwitz and Stieltjes calculus, with their ordinary
convergence hypotheses retained. The all-rank transport theorem
can be kept as a separate structural continuation. The Gamma basis
should follow the exact fixed-gap determinant reduction in the
harmonic-tail Gamma branch, with its inherited coordinate convention
printed explicitly.

The accompanying integration note lists these destinations, theorem
dependencies, a proposed status update, and the README patch.
Canonical acceptance should record the proved fixed-gap eventual
statement at its actual scope; it should not silently promote the
remaining full reconstruction or Gaussian conjectures.
